% Telugu chapter reader. English source at upstream/ remains pristine.
\documentclass[11pt,a4paper,oneside]{book}
\usepackage[margin=23mm,headheight=16pt]{geometry}
\usepackage{fontspec}
\setmainfont{Latin Modern Roman}
\setmathrm{Latin Modern Roman}
\newfontfamily{\telugufont}{NotoSerifTelugu-Regular.ttf}[Path=fonts/,Script=Telugu,Language=Telugu,BoldFont=NotoSerifTelugu-Bold.ttf,ItalicFont=NotoSerifTelugu-Regular.ttf,BoldItalicFont=NotoSerifTelugu-Bold.ttf]
\usepackage[Latin,Telugu]{ucharclasses}
\setTransitionsForLatin{\rmfamily}{}
\setTransitionTo{Telugu}{\telugufont}
\XeTeXgenerateactualtext=1
\usepackage{amsmath,amssymb,amsthm}
\usepackage{xparse,graphicx,xcolor,tikz}
\usepackage{float}
\floatplacement{figure}{H}
\usepackage[unicode,colorlinks=true,linkcolor=blue,urlcolor=blue]{hyperref}
\pdfstringdefDisableCommands{\def\telugufont{}\def\rmfamily{}}
\hypersetup{pdftitle={సమితులు - OpenLogic తెలుగు},pdfauthor={Open Logic Project; Telugu machine translation},pdfsubject={7 of 722 source units; first chapter tranche; provisional terminology}}
\setlength{\parindent}{0pt}
\setlength{\parskip}{6pt plus 1pt}
\linespread{1.18}
\tolerance=2000
\emergencystretch=3em
\raggedbottom
\renewcommand{\emph}[1]{{\bfseries #1}}
\renewcommand{\chaptername}{అధ్యాయం}
\renewcommand{\contentsname}{విషయ సూచిక}
\renewcommand{\figurename}{పటం}
\renewcommand{\proofname}{నిరూపణ}
\theoremstyle{definition}
\newtheorem{defn}{నిర్వచనం}[section]
\newtheorem{ex}[defn]{ఉదాహరణ}
\newtheorem{prop}[defn]{ప్రతిపాదన}
\newtheorem{thm}[defn]{సిద్ధాంతం}
\newtheorem{prob}{అభ్యాసం}[section]
\newenvironment{explain}{}{}
\newenvironment{digress}{\par}{\par}
\newenvironment{tagblock}[1]{}{}
\newcommand{\Setabs}[2]{\{#1:#2\}}
\newcommand{\Pow}[1]{\wp(#1)}
\newcommand{\Nat}{\mathbb{N}}
\newcommand{\Int}{\mathbb{Z}}
\newcommand{\Rat}{\mathbb{Q}}
\newcommand{\Real}{\mathbb{R}}
\newcommand{\Bin}{\mathbb{B}}
\newcommand{\PosInt}{\mathbb{Z}^{+}}
\newcommand{\lif}{\mathbin{\rightarrow}}
\newcommand{\tuple}[1]{\langle #1\rangle}
\newcommand{\len}[1]{\mathrm{len}(#1)}
\newcommand{\nicefrac}[2]{{#1}/{#2}}
\definecolor{oldiagcolorA}{rgb}{0.15,0.15,0.15}
\definecolor{oldiagcolorB}{rgb}{0.13,0.35,0.62}
\definecolor{oldiagcolorC}{rgb}{0.68,0.19,0.25}
\newcommand{\olasset}[1]{\centering\input{upstream/#1}}
\newcommand{\olfileid}[3]{\def\tepart{#1}\def\techapter{#2}\def\tesection{#3}}
\newcommand{\olchapter}[3]{\chapter{#3}}
\newcommand{\olsection}[1]{\section{#1}\label{\tepart:\techapter:\tesection:sec}}
\newcommand{\ollabel}[1]{\label{\tepart:\techapter:\tesection:#1}}
\NewDocumentCommand\olref{O{} O{} O{} m}{\if\relax\detokenize{#1}\relax\ref{\tepart:\techapter:\tesection:#4}\else\ref{#1:#2:#3:#4}\fi}
\makeatletter
\newcommand{\oliflabeldef}[3]{\@ifundefined{r@#1}{#3}{#2}}
\makeatother
\begin{document}
\frontmatter
\begin{titlepage}
\vspace*{25mm}
{\Huge\bfseries సమితులు\par}
\vspace{10mm}
{\Large OpenLogic తెలుగు అనువాదం\par}
\vspace{8mm}
మొదటి అధ్యాయ విడుదల: 722 మూల విభాగాలలో 7 విభాగాలు.
\par ఇది పూర్తి గ్రంథం కాదు. పూర్తి తెలుగు సంచికపై పని కొనసాగుతోంది.
\vfill
మూలం: Open Logic Project\par
యంత్ర అనువాదం; పదాల ఆధారాలు, తాత్కాలిక నిర్ణయాలు విడిగా నమోదయ్యాయి.
మానవ నిపుణుల సమీక్ష జరిగిందని ఈ సంచిక పేర్కొనదు.
\par గణిత విషయాలకు ఆంగ్ల మూలమే ఆధారం; తెలుగు సాక్ష్య గ్రంథాలు
భాషా వినియోగానికి ఆధారం మాత్రమే.
\vspace{8mm}
{\small మూల సంచిక: \texttt{9620cc73f9c8e0ad003c514a5d3748f29611c4c0}\par
\url{https://github.com/OpenLogicProject/OpenLogic}\par
\url{https://github.com/KokunoYumeto/OpenLogic-translations}}
\end{titlepage}
\section*{అనువాదం గురించి}
Open Logic Project మూలానికి Creative Commons Attribution 4.0 International
అనుమతి వర్తిస్తుంది. ఈ అనువాదం మూలంలో చేసిన మార్పుగా గుర్తించబడుతుంది.
మూల రచయితల వివరాలు, పూర్తి అనుమతి పత్రం, భాగాల ప్రత్యేక నిబంధనలు
సంపాదించగల మూలాలతోపాటు ఉంచబడ్డాయి.
\par ఈ అధ్యాయంలో ``మూలకాధారిత సమానత్వం'', ``ధర్మసంగ్రహం'' వంటి కొన్ని
పద నిర్ణయాలు తాత్కాలికమైనవి. వాటి కచ్చితమైన అర్థాన్ని ఇక్కడి
నిర్వచనాలే నిర్ణయిస్తాయి. సహజ సంఖ్యల సమితిలో $0$ను చేర్చే మూల
గ్రంథపు ఆనవాయితీని యథాతథంగా ఉంచాం.
\par తెలంగాణ, ఆంధ్ర ప్రదేశ్ సాక్ష్యాలను వేర్వేరుగా నమోదు చేశాం.
ఆ గ్రంథాలకు బహిరంగ పునర్ముద్రణ అనుమతి ఉందని భావించలేదు; వాటి పూర్తి
ప్రతులు ఈ విడుదలలో చేర్చబడవు.
\tableofcontents
\mainmatter
\input{build/sets-body.tex}
\end{document}
